% GENERATED FILE — DO NOT EDIT; authoritative prose is Markdown.
% source: theory/orthemic-core-formalization.md
% source-sha256: f55d98aa2dcaad231c769aa94e42fd0cec416361269c7d957b203669ca62dd72
% source: theory/orthemic-multi-actor-conflict-note.md
% source-sha256: aa7b1378b1f1983c3ec406fb2533a68c6242881c2d04bf34470632b37359d7b9
% source-qualification: research-stage-draft
% source-qualification: not-peer-reviewed
% source-qualification: not-externally-peer-reviewed
% source-qualification: not-empirically-validated
% source-qualification: candidate-terminology-not-adopted
\documentclass[10pt,letterpaper,twocolumn]{article}
\usepackage{amsmath}
\usepackage{amssymb}
\usepackage{booktabs}
\usepackage{geometry}
\usepackage{hyperref}
\usepackage{microtype}
\usepackage{natbib}
\usepackage{needspace}
\usepackage{xcolor}
\geometry{letterpaper,margin=0.75in}
\hypersetup{colorlinks=true,linkcolor=blue,urlcolor=blue,citecolor=blue}
\setlength{\emergencystretch}{3em}
\title{Orthemic Core Formalization — Orthing, Episodes, Metaorthemes, Pathway Adequacy}
\author{}
\date{}
\begin{document}
\twocolumn[
\maketitle

\begin{quote}
\textbf{Provenance.} Multi-model draft lineage with successive review passes; the original headers are preserved verbatim in \href{\detokenize{../docs/provenance/document-history.md}}{\texttt{docs/\allowbreak{}provenance/\allowbreak{}document-\allowbreak{}history.\allowbreak{}md}}, and model contributions in \href{\detokenize{../docs/provenance/model-contributions.md}}{\texttt{docs/\allowbreak{}provenance/\allowbreak{}model-\allowbreak{}contributions.\allowbreak{}md}}. Nothing here is experimentally validated.

\end{quote}
\textbf{Status (revision R6, 2026-07-20; fresh-session repository review completed; not external human peer review; not empirically validated):} canonical formalization of the orthing layer (episodes, metaorthemes, metaorthemmata, pathway adequacy), companion to the main manuscript's definitions. All claims are definitional/analytic unless marked empirical; nothing is experimentally validated — the deterministic fixture suite checks consistency only, and the Arm A/B/C terminology benchmark is designed but not run.

\textbf{Inherited notation} (normative registry: \href{\detokenize{../docs/notation-registry.yaml}}{\texttt{docs/\allowbreak{}notation-\allowbreak{}registry.\allowbreak{}yaml}}). $\mathcal{M}$, $\mathcal{O}$ = the universal occurrence and ortheme universes; a declared analysis $A$ activates a domain $\mathcal{M}_A \subseteq \mathcal{M}$ of orthemmata (concrete occurrences) and a repertoire $\mathcal{O}_A \subseteq \mathcal{O}$ of orthemes (repeatable operational state-types); under the standing single-analysis scope these are written $M$, $O$. $\operatorname{Inst}_A \subseteq \mathcal{M}_A \times \mathcal{O}_A$ = analysis-relative instantiation — the primitive; the manuscript's Definition 3 and §2.6 define the declared, versionable analysis $A$, whose components include the task $T = \operatorname{task}(A)$; $O^*(m; A)$ = actual profile (ground truth relative to $A$); $\Pi_A$ = the complete-profile space and $\Pi_A^\partial$ the partial-profile space (manuscript Definition 10); $\hat{p}_t(m) \in \Pi_A^\partial$ = inferred profile (belief; $\hat{O}$ survives only as the derived flattening of $\hat{p}$, §2.2); $\hat{C}_t(m)$ = candidate set of complete profiles; $\Omega$ = observation map; $\mu$ = a metaortheme; $\hat o$ = an inferred ortheme. \textbf{Abbreviation convention:} once a single analysis $A$ with $\operatorname{task}(A) = T$ is fixed by explicit declaration, task-subscripted forms ($\operatorname{Inst}_T$, $O^*_T(m)$, $\Pi_T$, $\mathcal{K}_T$, $\mathcal{R}_T$, $\mathcal{W}_T$) and the unsubscripted $M$, $O$ abbreviate their $A$-indexed counterparts. The shorthand is \textbf{forbidden} wherever more than one analysis is live — level-indexed audits (§1, §5.2), multi-actor episodes (§5.3), cross-version comparisons, execution-vs-review comparisons — where the full $O^*(m; A)$ form is required. (This block supersedes the earlier statement that treated the task-indexed relation as primitive; the retired symbols are listed in the notation registry's \texttt{retired\_symbols}.)


\bigskip\hrule\bigskip
\vspace{1em}
]

\section*{1. Type/token resolution (2A)}
\textbf{Orthing} (gloss: the doing of apprehension-and-handling) is a process TYPE: the rule-governed, evidence-updating operation kind by which an orthemma is individuated, observed, assigned a candidate and then inferred profile of orthemes, routed, validated, dispositioned, and revised under governing metaorthemes. It stands to its runs as "compilation" stands to "the 14:32 build of commit \texttt{abc123}." \textbf{An orthing episode} $e$ (gloss: one dated, situated run of that operation) is a concrete TOKEN: it has an actor, a time, an input occurrence, and a result, and two episodes are distinct occurrences even when inputs and outputs agree.

\textbf{Level-indexing.} Let $M^{(0)}$ be the base domain of orthemmata and $E^{(0)}$ the class of episodes placing members of $M^{(0)}$. When a higher audit asks "was this placement made correctly, on adequate evidence, under an adequate rule?", the episode is itself apprehended. This is an explicit \textbf{reification embedding}, not literal set inclusion — an episode is a different kind of thing from a jar of powder, and becomes a base occurrence of the higher audit only by being cast as one: for each level $n$,


\begin{verbatim}
ι_n : E^(n) ↪ M^(n+1)   (injective; ι_n(e) is "the episode e, considered
                         as a case to be audited")
\end{verbatim}
(gloss: yesterday's classification act, re-cast as today's case). Writing $\supseteq$ here would falsely assert episodes are already members of the base domain; the embedding records that a modeling step turns one into the other. The audit's orthemes at level $n+1$ are pathway state-types — \emph{evidence-sufficient placement}, \emph{stale-rule placement}, \emph{false closure} — precisely the §4 verdicts. This applies the manuscript's own licensed move ("metaorthemes can become first-order orthemes for a higher audit," §6) to the token side; the regress stops, as there, at a declared governance boundary.

\textbf{Why no separate "meta-orthemma" term.} The originating project notes float the noun "meta-orthemma" for the concrete operation. It is not adopted. Against the manuscript's own anti-vacuity standard, a new noun class must change at least one of representation, evidence, routing, or prediction; it changes none. \emph{Representation:} the embedding $\iota_n : E^{(n)} \hookrightarrow M^{(n+1)}$ already types the reified episode as an orthemma of the higher domain. \emph{Evidence:} what audits an episode is its signature record (§2) — a property of the domain, not of a new noun. \emph{Routing:} audits route episodes through the same lifecycle as any orthemma. \emph{Prediction:} nothing sayable with "meta-orthemma" is unsayable with "the reified episode $\iota_n(e) \in M^{(n+1)}$."

\textbf{The one honest counterargument, and its answer.} Episodes are unusual orthemmata: unlike a jar of powder, an episode carries a canonical, system-authored observation of itself — its trace. One might argue this structural difference deserves a name. Answer: the difference is real but is a \emph{subtype constraint on the domain} (episodes are orthemmata equipped with signatures), not a new pole of the type–token architecture. Register the subtype; do not coin the noun. If audits over episodes later prove to need operations generic placement lacks, revisit under the same test.

\textbf{Qualification (Decision 0002).} The rejection above concerns "meta-orthemma" as a name for the REIFIED EPISODE. Decision 0002 later adopted a different object in the same name-neighborhood: the \textbf{metaorthemma} $\bar{\mu}$ — an episode-local, case-bound CONFIGURATION TOKEN of a repeatable metaortheme (§2.2, §3, V3c in §4.1). It is not the whole episode and not the execution event; the rejection above stands for its original referent and does not adjudicate the token object (its revisit clause was exercised). The WORD "metaorthemma" remains a benchmark-gated candidate term; the OBJECT is adopted.


\bigskip\hrule\bigskip

\section*{2. Episode signature (2B)}

\subsection*{2.1 Candidate formulations compared}
\textbf{Candidate F (function style).} $e_\mu(m) = \hat o$: an episode governed by metaortheme $\mu$ maps orthemma $m$ to one ortheme; correct when $\hat o = o^*$. Virtues: minimal, mnemonic, correctness in one equation. Defects, each fatal for the paper's own commitments: (i) singleton output contradicts plural profiles (Def. 3) and route-sufficient apprehension (Def. 7); (ii) the episode is identified with its output — the function's \emph{value} — leaving no object to audit when the answer is right but the pathway defective; (iii) evidence, actor, time, route, residuals, and the successor occurrence have no slot; (iv) $\mu$ is one subscript, but real episodes run under several governing rules at once (§3.4).

\textbf{Candidate R (record/state-producing style).} The episode is a structured record — an element of a typed domain $E$ — with named components and projections; "output" is a designated sub-record. Virtues: audit-ready, plural by construction, separates run from result. Cost: heavier notation; needs a convention for the common quick case.

\textbf{Resolution:} adopt R as the definition; retain F as \emph{derived} notation for one placement inside an episode (2.3).


\subsection*{2.2 Recommended signature (formal)}
An \textbf{orthing episode} is a record


\begin{multline*}
e = \langle\, \mathrm{id};\ m, \kappa, v;\ x, H;\ \alpha, w, A, T, t;\\
\ \vec{\mu}, \mathrm{MetaTok}, \pi;\ \vec{C}, \hat{p};\ r;\\
\ \mathrm{estatus};\ \mathcal{Q};\ \delta;\ \mathrm{hand}_{\mathrm{in}},\\
 \mathrm{hand}_{\mathrm{out}};\ a, \mathrm{Succ} \,\rangle
\end{multline*}


with components (each glossed in ordinary language):


% breakable-row-table: data-rows=23 total-rendered-characters=5182 max-row-rendered-characters=1651
\par\medskip
\hrule height 0.8pt
\smallskip
% breakable-row: 1/23
\noindent\textbf{Component}: \texttt{id}\par
\noindent\textbf{Type}: episode identifier\par
\noindent\textbf{Gloss}: which run this is (referenced by handoffs and audits)\par
\smallskip
\hrule height 0.4pt
\smallskip
% breakable-row: 2/23
\noindent\textbf{Component}: $m$\par
\noindent\textbf{Type}: $M$\par
\noindent\textbf{Gloss}: the input orthemma — the concrete case handled\par
\smallskip
\hrule height 0.4pt
\smallskip
% breakable-row: 3/23
\noindent\textbf{Component}: $\kappa = \operatorname{id}(m)$\par
\noindent\textbf{Type}: identity key\par
\noindent\textbf{Gloss}: which thing this is (survives change)\par
\smallskip
\hrule height 0.4pt
\smallskip
% breakable-row: 4/23
\noindent\textbf{Component}: $v = \operatorname{ver}(m)$\par
\noindent\textbf{Type}: version\par
\noindent\textbf{Gloss}: which state/edition of that thing\par
\smallskip
\hrule height 0.4pt
\smallskip
% breakable-row: 5/23
\noindent\textbf{Component}: $x = \Omega_e(m)$\par
\noindent\textbf{Type}: observation\par
\noindent\textbf{Gloss}: what the episode actually saw\par
\smallskip
\hrule height 0.4pt
\smallskip
% breakable-row: 6/23
\noindent\textbf{Component}: $H$\par
\noindent\textbf{Type}: ordered typed evidence sequence\par
\noindent\textbf{Gloss}: the evidence gathered in order, each item with property class, scope, provenance, and validity/expiry\par
\smallskip
\hrule height 0.4pt
\smallskip
% breakable-row: 7/23
\noindent\textbf{Component}: $\alpha$\par
\noindent\textbf{Type}: actor\par
\noindent\textbf{Gloss}: who/what executed (human, validator, model, pipeline, institution — populated even for mechanical executors)\par
\smallskip
\hrule height 0.4pt
\smallskip
% breakable-row: 8/23
\noindent\textbf{Component}: $w$\par
\noindent\textbf{Type}: episode warrant\par
\noindent\textbf{Gloss}: who/what authorized acting, with scope (the \#12 hook); distinct from evidence and from executor identity\par
\smallskip
\hrule height 0.4pt
\smallskip
% breakable-row: 9/23
\noindent\textbf{Component}: $A$\par
\noindent\textbf{Type}: declared analysis (identifier + version)\par
\noindent\textbf{Gloss}: the analysis the placement is relative to (manuscript §2.6); result correctness (V1) is judged against $O^*(m; A)$; within the episode's scope the task-subscripted spaces below abbreviate their $A$-indexed forms\par
\smallskip
\hrule height 0.4pt
\smallskip
% breakable-row: 10/23
\noindent\textbf{Component}: $T = \operatorname{task}(A)$\par
\noindent\textbf{Type}: task\par
\noindent\textbf{Gloss}: the task component of $A$, retained as a separate readable component\par
\smallskip
\hrule height 0.4pt
\smallskip
% breakable-row: 11/23
\noindent\textbf{Component}: $t$\par
\noindent\textbf{Type}: time interval\par
\noindent\textbf{Gloss}: when the episode ran\par
\smallskip
\hrule height 0.4pt
\smallskip
% breakable-row: 12/23
\noindent\textbf{Component}: $\vec{\mu} = (\mu_1, \dots, \mu_k; \preceq)$\par
\noindent\textbf{Type}: metaorthemic configuration\par
\noindent\textbf{Gloss}: the governing rules in force, with precedence order $\preceq$\par
\smallskip
\hrule height 0.4pt
\smallskip
% breakable-row: 13/23
\noindent\textbf{Component}: $\operatorname{MetaTok}(e) = \{\bar\mu_1, \dots, \bar\mu_j\}$\par
\noindent\textbf{Type}: concrete metaorthemmata (Decision 0002)\par
\noindent\textbf{Gloss}: the case-bound CONFIGURATION TOKENS of governing types actually applied. Each $\bar{\mu}$ records/references: its own identity and lineage; $(\mu, \operatorname{ver}(\mu))$ via the internal typing $\operatorname{MetaInst}(\bar{\mu}, \mu)$; $(A(e), \operatorname{ver}(A(e)))$ with $\operatorname{Compatible}(\bar{\mu}, A(e))$ — the token binds case-specific values WITHIN the declared analysis, references any value $A$ fixes uniquely, and never overrides $A$ without an explicit new analysis version; the target $(\kappa, v)$; governed component $g$; the case-specific binding map $B$ (reference frame, tolerance value, fixture, success surface, …); scope $\sigma$ incl. the claims in $\mathcal{Q}$ that depend on it; references to policy, evidence selector, instrument/tool and calibration provenance; the \textbf{binder} (actor/process that made the binding) with binding warrant $w_{\mathrm{bind}}$, kept DISTINCT from the designated executor; binding time $t_{\mathrm{bind}}$; validity/expiry/supersession. The token REFERENCES the episode's evidence $H$, trace, and output — it never absorbs them (the application-event view $\operatorname{ApplyEvent}(\bar{\mu}, e) = \langle\bar{\mu}, \operatorname{Trace}_e |_{\bar{\mu}}\rangle$ is derived, not primitive). \textbf{Zero-burden rule:} a configuration with no material case-specific binding, non-default scope, instrument/calibration, or independent validity condition gets no explicit token, and $\operatorname{V3c} \notin \operatorname{ReqPath}(e)$ (status recorded not-applicable, with reason)\par
\smallskip
\hrule height 0.4pt
\smallskip
% breakable-row: 14/23
\noindent\textbf{Component}: $\pi$\par
\noindent\textbf{Type}: policy\par
\noindent\textbf{Gloss}: the concrete procedure executed under $\vec{\mu}$\par
\smallskip
\hrule height 0.4pt
\smallskip
% breakable-row: 15/23
\noindent\textbf{Component}: $\vec{C}$\par
\noindent\textbf{Type}: typed candidate families\par
\noindent\textbf{Gloss}: open alternatives per uncertainty axis: $C^{\mathrm{id}} \subseteq M$ (which occurrence), $C^{\mathrm{profile}} \subseteq \Pi_A$ (which profile), $C^{\mathrm{cause}} \subseteq \mathcal{K}_A$ (which cause), $C^{\mathrm{route}} \subseteq \mathcal{R}_A$ (which route), $C^{\mathrm{warrant}} \subseteq \mathcal{W}_A$ (which warrant state) — competing hypotheses may themselves be PROFILES, so candidates are not forced into single orthemes\par
\smallskip
\hrule height 0.4pt
\smallskip
% breakable-row: 16/23
\noindent\textbf{Component}: $\hat{p} \in \Pi_A^\partial$\par
\noindent\textbf{Type}: inferred placement\par
\noindent\textbf{Gloss}: \textbf{exactly one partial profile} — the profile actually placed, determined on resolved axes and open on the rest; \textbf{never a set of complete profiles} (live alternatives belong in $\hat{C} \subseteq \Pi_A$, and optional weights in the separate belief layer) and \textbf{never coerced to a singleton ortheme}; $\hat{O} \subseteq O$ survives only as the derived flattening of $\hat{p}$. (R4 repair, Decision 0012: the pre-R4 gloss "or a bounded set of profiles if unresolved" typed one symbol into two spaces and is withdrawn; archived patches retain it as history.)\par
\smallskip
\hrule height 0.4pt
\smallskip
% breakable-row: 17/23
\noindent\textbf{Component}: $r$\par
\noindent\textbf{Type}: route\par
\noindent\textbf{Gloss}: the operation/owner the case was sent to\par
\smallskip
\hrule height 0.4pt
\smallskip
% breakable-row: 18/23
\noindent\textbf{Component}: $\operatorname{estatus}$\par
\noindent\textbf{Type}: evidence status map\par
\noindent\textbf{Gloss}: per placed claim: validated / provisional / stale / absent\par
\smallskip
\hrule height 0.4pt
\smallskip
% breakable-row: 19/23
\noindent\textbf{Component}: $\mathcal{Q}$\par
\noindent\textbf{Type}: claim ledger\par
\noindent\textbf{Gloss}: per claim: proposition; target occurrence/version; property class (\#3); required success surface (\#14); evidence IDs; warrant where relevant; verification status; explicit non-claims\par
\smallskip
\hrule height 0.4pt
\smallskip
% breakable-row: 20/23
\noindent\textbf{Component}: $\delta$\par
\noindent\textbf{Type}: residual disposition map\par
\noindent\textbf{Gloss}: per obligation: unresolved / deferred / transferred / owner-assigned / risk-accepted / validated-resolved (issue \#6's six values). Claims ($\mathcal{Q}$) are assertions; residuals ($\delta$) are obligations — never merged\par
\smallskip
\hrule height 0.4pt
\smallskip
% breakable-row: 21/23
\noindent\textbf{Component}: $\operatorname{hand\_in}, \operatorname{hand\_out}$\par
\noindent\textbf{Type}: handoff records\par
\noindent\textbf{Gloss}: incoming/outgoing packets: packet id+version, sender/receiver episode ids, subject identity, state claims, evidence references, authority claims (the \#15 hook)\par
\smallskip
\hrule height 0.4pt
\smallskip
% breakable-row: 22/23
\noindent\textbf{Component}: $a$\par
\noindent\textbf{Type}: action\par
\noindent\textbf{Gloss}: what was done to or with the case\par
\smallskip
\hrule height 0.4pt
\smallskip
% breakable-row: 23/23
\noindent\textbf{Component}: $\operatorname{Succ} \subseteq M$\par
\noindent\textbf{Type}: successor set\par
\noindent\textbf{Gloss}: the occurrences the action created/affected, as LABELED lineage edges (which action produced which successor) — zero, one, or many; the earlier forced single $m' = \operatorname{succ}(m,a)$ is corrected\par
\smallskip
\hrule height 0.8pt
\par\medskip
$\Pi_A$ — with $\Pi_T$ its licensed shorthand under the standing convention — is the \textbf{profile space} of manuscript Definition 10, R3 general form: $\Pi_A \subseteq \mathcal{P}(\mathcal{O}_A)$, the set of admissible complete profiles under the constraints declared by $A$. Factorized axes are one \emph{permitted representation family}, not a universal ontology: where declared, a complete profile selects exactly one admissible value per \textbf{applicable} alternatives-marked axis, records an explicit \texttt{not-applicable} value on an objectively inapplicable axis (a fact about the case, never uncertainty), and contains any declared-admissible — possibly empty — subset per co-holding axis, consistent with the declared cross-axis constraints. Objective absence, objective inapplicability, epistemic openness, evidence absence, and candidate plurality are five distinct states; only the first two live in profiles. $\Pi_A^\partial$ is the partial-profile space (resolutions may be left undetermined), with $\Pi_A \subseteq \Pi_A^\partial$ (gloss: every complete way the case could be placed, and every honest partial state of knowledge about it). This and the labeled successor edges close the arrow the manuscript lacks from action back into $M$: placements verified for $(\kappa, v)$ do not transport to $(\kappa, v')$ without a lineage argument.

\textbf{Plain-language key for governed-component letters used below:} $O$ = the ortheme repertoire (which state-types exist); $I$ = individuation/identity (what counts as the same case, versioning); $E$ = evidence procedure (what evidence counts, at what grade); $D$ = disclosure/uncertainty handling (what may be left open and how it is declared); $R$ = routing (where cases are sent); $V$ = validation/closure (what may be called done, at what standard); $W$ = warrant/authorization-state classification (which warrant states exist — authorized vs factually established — and what each licenses).


\subsection*{2.2a Endpoint plus bounded trace}
The record above is the REQUIRED endpoint. Because orthing is defined as evidence-updating, an endpoint alone cannot always adjudicate pathway verdicts; where those verdicts (V2b-P truth-conduciveness, V3d executor fidelity, V6 robustness — §4) must be adjudicated from the episode's own record, an ordered \textbf{trace} is also required:


\begin{verbatim}
Trace_e = (s_0, u_1, s_1, …, u_k, s_k)
\end{verbatim}
where each state $s_i$ snapshots the current evidence, candidate families, inferred placement, route, and open residuals, and each update $u_i$ records an evidence acquisition, selection, routing, action, or revision step (with the rule edition $\operatorname{ver}(\mu)$ in force). \textbf{Governance declares the trace granularity}; trivial cases carry no trace, and maximal per-event logging is never a blanket requirement.

\textbf{Partial applicability.} Not every episode has a route, a mutation, a closure claim, or open residuals. Components are defined \textbf{where applicable}: a read-only classification episode has $a = \bot$ (no action), $\operatorname{Succ} = \emptyset$ (no successor created), possibly $r = \bot$ (no downstream route), and an empty $\delta$. Where a component is $\bot$, the verdicts that quantify over it (V4a route, V5 closure — §4) are \textbf{not asserted of that episode} rather than counted as failures; formally, pathway adequacy conjoins over the governance-DERIVED required set $\operatorname{ReqPath}(e)$ (§4.1, Decision 0003) — applicability is derived from the declared analysis, episode shape, risk class, claims, and governance, never a discretionary post-hoc list, and every verdict excluded from $\operatorname{ReqPath}(e)$ carries a recorded \texttt{not-applicable} reason on the episode record. So "no route" (deliberately none, with reason) is distinguished from "route omitted" (a defect), and "not tested" is never "not applicable".


\subsection*{2.3 Three separations, stated once}

\begin{enumerate}
\item \textbf{Episode ≠ output.} $\operatorname{out}(e) = \langle \hat{p}, r, \operatorname{estatus}, \mathcal{Q}, \delta, a \rangle$ is a projection of $e$; episodes with identical outputs can differ in pathway and hence in every §4 verdict.

\item \textbf{Episode ≠ policy.} $\pi$ is repeatable and undated; $e$ is one dated execution of $\pi$. Reliability claims attach to $\pi$ (and $\vec{\mu}$); correctness claims attach to $e$.

\item \textbf{Placement inside an episode ≠ the episode.} The compact function style becomes the derived judgment

\end{enumerate}

\begin{multline*}
e \models_{\mu} (m : \hat{o}) \iff\\
 \hat{o} \in \hat{O}(e) \ \text{and}\ \mu \in \vec{\mu}(e)\ \text{governed that placement}
\end{multline*}


(gloss: "in episode $e$, under rule $\mu$, the case was placed under $\hat o$"). One placement is correct iff $\hat o \in O^*(m; A(e))$; the placed profile is correct at the governed level iff $\hat{p}(e)$ agrees with $O^*(m; A(e))$ restricted to that level, or the weaker route-sufficient variant (§4) — $A(e)$ being the analysis recorded in the episode (§2.2).


\bigskip\hrule\bigskip

\section*{3. Metaortheme semantics (2C)}
\textbf{Type vs token (Decision 0002).} This section describes the REPEATABLE side: the metaortheme type $\mu$ and its paired meta-policy $\pi_\mu$. Their concrete, case-bound instantiation inside an episode is the \textbf{metaorthemma} $\bar{\mu}$ (§2.2), typed by the internal relation $\operatorname{MetaInst}(\bar{\mu}, \mu)$ — a typing of tokens, not a second ground-truth primitive — and subordinated to the episode's analysis by $\operatorname{Compatible}(\bar{\mu}, A(e))$. Four layers stay separate: type $\mu$; policy $\pi_\mu$; token $\bar{\mu}$ (binding); execution of the token (episode trace, judged by V3d). The corresponding adequacy questions are independent: adequate type, adequate policy, correctly bound token, faithful execution.


\subsection*{3.1 Three models}

\begin{itemize}
\item \textbf{(A) Higher-order STATE/TYPE.} A metaortheme is an ortheme \emph{of the governing context}: the context instantiates one of several competing higher-order states (e.g., evidence is provenance-grade or appearance-grade), and that state governs.

\item \textbf{(B) Higher-order handling RULE.} A metaortheme is a repeatable rule over the machinery: "never place on appearance alone," "quarantine on stale version." The rule \emph{is} the object.

\item \textbf{(C) Split model.} A metaortheme is a \textbf{metaorthemic configuration} $\mu = \langle g, S_\mu, \operatorname{select}_\mu, \operatorname{prov}(\mu), \operatorname{ver}(\mu)\rangle$ — governed component $g$, declared competing higher-order states $S_\mu$, selecting-evidence procedure $\operatorname{select}_\mu$ that determines which state obtains, provenance, version — together with a distinct \textbf{meta-policy} $\pi_\mu$ prescribing conduct conditional on the state (gloss: the \emph{distinction the rule consults} and the \emph{rule that consults it} are two objects, as Def. 9 already insists in prose).
\textbf{Governed-component boundary (adjudicated; one boundary for all documents).} $g \in \{\mathrm{O}, \mathrm{I}, \mathrm{E}, \mathrm{D}, \mathrm{R}, \mathrm{V}, \mathrm{W}\}$ — the mapping/handling components plus warrant/authorization-state classification (plain- language key in §2.2; $W$ is required by worked example 5). Excluded as governed components: \textbf{objectives/loss} (the Grok exchange remains the standing negative example of an objective mislabeled a metaortheme); the \textbf{task} $T$ (changing the task changes the problem — rules are task-\emph{indexed}, never task-\emph{governing}); and the \textbf{governance meta-level itself} (the precedence order $\preceq$ and revision rights are parameters fixed at the declared governance boundary, where the regress terminates). The paper memo's earlier Def-9 "transformation" framing is reconciled as: a metaortheme \emph{change} induces a transformation of the mapping subsystem — the transformation is the effect of revising $\mu$, not $\mu$ itself.

\end{itemize}

\subsection*{3.2 Minimum-specificity test against the five worked examples}

% breakable-row-table: data-rows=5 total-rendered-characters=1363 max-row-rendered-characters=287
\par\medskip
\hrule height 0.8pt
\smallskip
% breakable-row: 1/5
\noindent\textbf{Owner example}: appearance-only vs provenance/discriminating evidence\par
\noindent\textbf{Governed}: $E$ (evidence procedure)\par
\noindent\textbf{Competing states}: \{appearance-grade, provenance-grade\}\par
\noindent\textbf{Selecting evidence}: source record or discriminating test\par
\noindent\textbf{Counterfactual difference}: appearance-grade placements are reversible on test; provenance-grade support closure\par
\noindent\textbf{Policy relation}: policy chooses which test to run before placing\par
\smallskip
\hrule height 0.4pt
\smallskip
% breakable-row: 2/5
\noindent\textbf{Owner example}: current vs stale version\par
\noindent\textbf{Governed}: identity/versioning on $M$\par
\noindent\textbf{Competing states}: \{current, stale\}\par
\noindent\textbf{Selecting evidence}: lineage check $\operatorname{ver}(m)$ vs latest\par
\noindent\textbf{Counterfactual difference}: verdicts on stale versions do not transport\par
\noindent\textbf{Policy relation}: policy quarantines or re-derives on stale\par
\smallskip
\hrule height 0.4pt
\smallskip
% breakable-row: 3/5
\noindent\textbf{Owner example}: route-sufficient vs identity-complete\par
\noindent\textbf{Governed}: $R$,$V$ (routing, validation)\par
\noindent\textbf{Competing states}: \{route-sufficient, identity-complete\}\par
\noindent\textbf{Selecting evidence}: task declaration + risk class\par
\noindent\textbf{Counterfactual difference}: act-now vs investigate-further diverge\par
\noindent\textbf{Policy relation}: policy releases the route while holding residuals open\par
\smallskip
\hrule height 0.4pt
\smallskip
% breakable-row: 4/5
\noindent\textbf{Owner example}: closure-safe vs false closure\par
\noindent\textbf{Governed}: $V$ (closure)\par
\noindent\textbf{Competing states}: per-burden \{resolved, deferred, transferred, risk-accepted\} vs "all done"\par
\noindent\textbf{Selecting evidence}: residual ledger $\delta$\par
\noindent\textbf{Counterfactual difference}: falsely closed cases silently accrue debt; safe closure is reopenable by record\par
\noindent\textbf{Policy relation}: policy may claim completion only at the ledger's level\par
\smallskip
\hrule height 0.4pt
\smallskip
% breakable-row: 5/5
\noindent\textbf{Owner example}: authorized vs factually established\par
\noindent\textbf{Governed}: $W$ (warrant-state classification)\par
\noindent\textbf{Competing states}: \{authorized, established, both, neither\}\par
\noindent\textbf{Selecting evidence}: authorization record vs validating evidence\par
\noindent\textbf{Counterfactual difference}: authorized-but-unestablished placements must carry a revisit obligation\par
\noindent\textbf{Policy relation}: policy tags warrant type and never launders one into the other\par
\smallskip
\hrule height 0.8pt
\par\medskip
Verdict per model: \textbf{A} captures competing states but has no slot for the last column — a state with no conduct attached governs nothing. \textbf{B} captures conduct but fails the manuscript's anti-vacuity clause (iv) — a rule renamed as a noun — and cannot say what two rules consulting the \emph{same} distinction (quarantine vs re-derive on staleness) have in common. \textbf{C} fills every cell of every row; the state/policy split is load-bearing exactly where B collapses (same $S_\mu$, different $\pi_\mu$).


\subsection*{3.3 Recommendation}
\textbf{Adopt (C).} It is the only model under which all five worked examples pass minimum specificity; it matches Definition 9's prose; it makes rule-adequacy (§4, V3a) a predicate on $\mu$ separable from executor error in $\pi_\mu$; and it gives lesson-lift a target (revise $S_\mu$ or $\operatorname{select}_\mu$) distinct from retraining conduct (revise $\pi_\mu$).

\textbf{Preserved downsides of (C):} two objects per rule is real bookkeeping; most working systems will implement (B) and only reconstruct the split when audited — (C) is the audit-grade description, not a claim about how practitioners write rules; boundary cases blur state and policy (a numeric threshold is at once a state boundary and a conduct trigger); and the split invites proliferation — every \texttt{if}-statement can be dressed as configuration-plus-policy — which the safeguards below must block.


\subsection*{3.4 Plurality, conflict, precedence, provenance, revision, safeguards}

\begin{itemize}
\item \textbf{Several metaorthemes per episode.} $\vec{\mu}(e) = (\{\mu_1, \dots, \mu_k\}, \preceq)$. Configurations governing \emph{disjoint} components compose freely (their constraints conjoin); a typical episode runs under a version rule, an evidence-grade rule, and a closure rule at once.

\item \textbf{Conflict and precedence.} $\mu_i, \mu_j$ conflict at $e$ when their meta-policies prescribe incompatible constraints on the same component under the states that actually obtain. Resolution only via the declared strict partial order $\preceq$ (gloss: which rule outranks which, fixed by governance in advance). A conflict unresolved by $\preceq$ is an ANDON condition — stop or escalate; silent override is itself a metaorthemic error.

\item \textbf{Provenance of the governing rule.} $\operatorname{prov}(\mu) =\allowbreak{} \langle \mathrm{authority},\allowbreak{} \mathrm{warrant},\allowbreak{} \mathrm{scope},\allowbreak{} \operatorname{ver}(\mu) \rangle$. A rule of unknown provenance is a stale-evidence problem one level up; episodes record $\operatorname{ver}(\mu)$ so audits can scope which placements ran under which edition. Governing \emph{instructions} recovered from records obey the same discipline: existence, authenticity, recovery, authorization, and current force are five distinct predicates, and an authentic recovered directive may be stale as governance — the worked stale-steer case, with its directive record and fixture F6, is \texttt{examples/\allowbreak{}compaction-\allowbreak{}stale-\allowbreak{}steer.\allowbreak{}md} (Decision 0006).

\item \textbf{Revision after lesson-lift (impact-scoped).} $\rho : (\mu, \mathrm{lesson}) \mapsto \mu'$ increments $\operatorname{ver}(\mu)$ and computes an \textbf{impact function} $\operatorname{affected}(\mu, \mu') \to$ the set of dependent claim/episode scopes: exactly the prior episodes whose placement CONSULTED the revised state-distinction or selector where the two editions disagree. Only those episodes' $\epsilon$ is downgraded from validated to provisional pending re-check (gloss: when the rule was wrong, the wins that depended on the wrong part are reopened — not every episode that ever ran under the edition). Blanket reopening of the whole edition cohort is corrected to this scoped form.

\item \textbf{State families need not partition, and selection may be undetermined.} $S_\mu$ declares the competing higher-order states, but nothing forces them to be mutually exclusive or jointly exhaustive unless the configuration declares them so (an exclusivity marking, as for profile axes); and $\operatorname{select}_\mu$ is a partial procedure — on a given episode it may return one state, several co-holding states, or \textbf{undetermined}. An undetermined selection is handled like any other unresolved placement: the meta-policy must declare conduct for it (stop, escalate, default conservatively), and silently assuming a partition where none was declared is a configuration defect (V3a).

\item \textbf{Safeguards against inflation.} Admit a candidate $\mu$ only if: (i) $g$ is named; (ii) $S_\mu$ is declared \emph{in advance}, not read off one incident; (iii) switching the obtaining state changes validated placement, risk, or constraints in at least one episode class; (iv) it is not a first-order ortheme or a tunable parameter of $\pi$ in disguise. This is the manuscript's anti-vacuity test with (ii)–(iii) sharpened to episode classes. A good policy that consults no in-advance-declared competing higher-order states is just a good policy. \textbf{Concept admission vs word utility are separate tests:} the four conditions above admit the DISTINCTION; whether the \emph{word} "metaortheme" (or any coinage) earns keep over ordinary words is decided ONLY by the terminology benchmark (the terminology evaluation protocol under \texttt{terminology/\allowbreak{}}) — a real distinction may be admitted while its word is retired. (The former safeguard (v), which mixed the two tests, is deleted.)

\end{itemize}

\bigskip\hrule\bigskip

\section*{4. Pathway adequacy (2D)}

\subsection*{4.1 Separate verdicts}
All verdicts are predicates on an episode $e$ (some claim-indexed). \textbf{The verdicts diagnose distinct dimensions and are not identified with one another; they are NOT assumed pairwise logically independent.} Every definitional or logical implication is declared explicitly in the implication table at the end of this section; no implication may be inferred merely from proximity in the verdict numbering (Decision 0003; this supersedes the earlier "none entails another"). Verdict labels follow the normative registry (\href{\detokenize{../docs/verdict-registry.yaml}}{\texttt{docs/\allowbreak{}verdict-\allowbreak{}registry.\allowbreak{}yaml}}, Decision 0004): semantic IDs are authoritative in machine-readable records, and the display aliases below are the prose forms. \textbf{Compact core (main text):} result correctness; evidential support; procedure truth-conduciveness; rule adequacy; executor fidelity; route safety; closure truthfulness; robustness. \textbf{Full vector (this table)} — the finer splits matter to audits; ordinary discussion may use the compact core:


\begin{itemize}
\item \textbf{V1 — placement/profile correctness.} $\operatorname{correct}(e)$: $\hat{p}(e)$ agrees with $O^*(m; A(e))$ at the governed level; weaker $\operatorname{route-correct}(e)$: $\hat{p}(e)$ is route-sufficient and every placed claim is in $O^*(m; A(e))$ (gloss: the answer was right, for enough of the profile to act on — judged against the analysis the episode itself records). \textbf{Aggregation, stated once:} V1 is claim-wise at bottom — $\operatorname{V1}_q(e)$ holds iff placed claim $q$ is true of $(m; A(e))$ — and the episode-level verdict aggregates by CONJUNCTION over the claims placed at the governed level (for $\operatorname{correct}(e)$, additionally requiring that the placed profile is determined at that level); the route-sufficient reading conjoins over the placed claims only. No averaging or majority rule is admissible for V1. \textbf{Result-side: V1 is never a conjunct of $\operatorname{PathwayAdequate}$.}

\item \textbf{V2a — evidential support.} The typed evidence $H$, within its declared scopes, meets the DECLARED standard for each placed claim (gloss: enough of the right kind of evidence, by the rulebook's own bar).

\item \textbf{V2b-P — configured-procedure truth-conduciveness (pathway-side; Decision 0003).} NON-FACTIVE. The procedure family actually instantiated in this episode — under its governing configuration $\vec{\mu}$, applicable metaorthemmata $\operatorname{MetaTok}(e)$, and execution mode — satisfies the predeclared reliability criterion over its DECLARED reference class. This is not an unrestricted claim about the abstract procedure type. For each placed claim $q$, the reliability specification records or references $\operatorname{RelSpec}_q(e) =\allowbreak{} \langle \text{declared reference class};\allowbreak{} \text{relevant case/risk stratum};\allowbreak{} \text{reliability metric};\allowbreak{} \text{threshold or tolerance};\allowbreak{} \text{perturbation or comparison }\allowbreak{}\text{family where applicable};\allowbreak{} \text{evaluation protocol};\allowbreak{} \text{evidence used to establish reliability};\allowbreak{} \text{version and validity conditions} \rangle$. The reference class, criterion, and threshold must be fixed independently of this episode's outcome — never selected after seeing the result merely to rescue or condemn it. The current episode's correctness may contribute to LATER aggregate reliability evidence, but one current correct result cannot by itself establish V2b-P. \textbf{$\operatorname{V2b-P}(e)$ does not entail $\operatorname{V1}(e)$:} a reliable configured procedure may produce a rare error. \emph{Default criterion:} a \textbf{sensitivity}-style counterfactual (the check passes because claims of this class are true, not alongside them); soundness, completeness, calibration, and perturbation-stability variants are admissible declarations, and fixtures may test each.

\item \textbf{V2b-T — token-level truth linkage (result-side audit annotation; EXCLUDED from the pathway core).} FACTIVE and claim-wise: $\operatorname{V2b-T}_q(e)$ holds when THIS placed claim $q$ was correct through the truth-relevant evidential mechanism rather than merely correct alongside it; by definition $\operatorname{V2b-T}_q(e) \to \operatorname{V1}_q(e)$ (truth of that claim). One claim-level $\operatorname{V2b-T}_q$ does not by itself entail correctness of the whole profile — a profile-level reading requires every placed claim covered and an explicit aggregation rule. It remains reportable for stopped-clock/Gettier diagnoses; its factivity is exactly why it sits outside $\operatorname{PathwayAdequate}$ (including it would re-import result correctness). Token-local PROCEDURAL defects are not represented by V2b-T alone: non-factive defects must appear under the appropriate pathway verdicts (V2a, V2b-P, V2c, V3a, V3b, V3c, V3d, or V6).

\item \textbf{V2c — evidence currentness/provenance.} Each load-bearing evidence item is current for $(\kappa, v)$ and carries admissible provenance (gloss: not stale, not unsourced — separated from V2a because evidence can meet the standard's \emph{strength} while being the wrong vintage).

\item \textbf{V3a — governing-configuration adequacy.} $\operatorname{adeq}(\vec{\mu}, e)$: each $\mu_i$ in force is adequate for the case's risk class — $S_\mu$ separates the higher-order states this case class can occupy and $\operatorname{select}_\mu$ can actually discriminate them (gloss: the rulebook was strong enough, whoever executed it).

\item \textbf{V3b — meta-policy adequacy.} The meta-policy $\pi_\mu$/procedure $\pi$ is well-formed for the configuration it runs under (a sound rulebook can carry an ill-formed procedure).

\item \textbf{V3c — governing-token adequacy (Decision 0002).} $\operatorname{V3c}(e) \Leftrightarrow\allowbreak{} \wedge_{\bar{\mu} \in \operatorname{MetaTok}(e)} \operatorname{TokenAdequate}(\bar{\mu}, e)$, where $\operatorname{TokenAdequate}(\bar{\mu}, e) \Leftrightarrow\allowbreak{} \operatorname{MetaInst}(\bar{\mu}, \mu) \wedge\allowbreak{} \operatorname{Compatible}(\bar{\mu}, A(e)) \wedge\allowbreak{} \operatorname{Anchored}(\bar{\mu}, \kappa(e), v(e)) \wedge\allowbreak{} \operatorname{ScopeCorrect}(\bar{\mu}, \mathcal{Q}(e)) \wedge\allowbreak{} \operatorname{Current}(\bar{\mu}, t(e)) \wedge\allowbreak{} \operatorname{Provenanced}(\bar{\mu}) \wedge\allowbreak{} \operatorname{AuthorizedBinding}(\bar{\mu})$ — every applicable concrete metaorthemma was correctly instantiated, analysis-compatible, occurrence-anchored, correctly scoped to the claims it served, current, provenanced, and bound under authority. PER-TOKEN statuses are preserved alongside the episode-level conjunction (the auditor must see WHICH governing application failed, e.g. \texttt{μ̄\_2: stale calibration; μ̄\_3: wrong claim scope}). Isolates the failure mode the vector previously could not name: correct standard (V3a) + sound policy (V3b) + faithful execution (V3d) + \textbf{defective case-specific binding} — wrong reference plane, wrong-role tolerance, expired calibration, wrong fixture or success surface. $\operatorname{MetaTok}(e) = \emptyset \Rightarrow \operatorname{V3c} \notin \operatorname{ReqPath}(e)$ (zero-burden rule; status recorded not-applicable, with reason). The Decision-0004 lettering matches the conceptual adequacy chain: V3a (configuration) → V3b (policy) → \textbf{V3c (binding)} → V3d (execution) → V3e (decision-time).

\item \textbf{V3d — executor fidelity.} The actor $\alpha$ actually executed $\pi$ as written. An adequate rule + adequate policy + infidelic executor is a distinct failure mode routing to a different remedy (retrain/replace the executor, not rewrite the rule).

\item \textbf{V3e — ex-ante procedural justification.} At decision time $t$, given exactly the evidence then available, the placement was the reasonable one under the declared loss/decision standard (gloss: blameless-at-the- time — indexed to decision time, unlike V2a/V2b-P/V2b-T which are audit-time verdicts about support and truth-connection).

\item \textbf{V4a — route safety/admissibility.} $r$ is admissible under the constraints in force; correct-route-unavailable is recorded as routing failure, not diagnostic uncertainty.

\item \textbf{V4b — route near-optimality.} $r$ is near-optimal given $\hat{p}(e)$ and the constraints (gloss: safe ≠ best; a safe but wasteful route fails V4b only).

\item \textbf{V5 — closure truthfulness.} Every residual in $\delta$ has an admissible disposition with traceable ownership, and the completion claim matches the ledger — no collapse of \{deferred, transferred, risk-accepted\} into "resolved."

\item \textbf{V6 — robustness under neighboring perturbations.} Robustness is always relative to a declared \textbf{perturbation specification} $\operatorname{PerturbSpec}(e) =\allowbreak{} \langle \text{varied fields};\allowbreak{} \text{invariant fields};\allowbreak{} \text{generator or enumeration};\allowbreak{} \text{size or measure};\allowbreak{} \text{tolerance} \rangle$: which episode/input fields are deliberately varied (input version bumped, marker string removed or spoofed, evidence reordered or withheld, near-identical sibling substituted), which are held invariant (the policy $\pi$, the governing configuration $\vec{\mu}$, the analysis $A$ and its version), and how the neighborhood is generated. This induces the perturbation relation $P \subseteq E \times E$ and neighborhood $N(e) = \{e' : (e, e') \in P\}$ (gloss: the same machine on the cases next door). $\operatorname{robust}(e)$ holds iff EITHER (i) the declared neighborhood is finite and enumerated and the empirical proportion of V1 failures over it is at or below the declared tolerance, OR (ii) the specification declares a probability measure over the perturbation family and the failure probability under that measure is at or below tolerance — an undeclared "failure rate" over an unspecified infinite neighborhood is not a well-formed V6 claim. \textbf{Metaorthemma rebinding under perturbation:} where a perturbation bumps the input version, each governing token $\bar{\mu}$ is re-bound to the perturbed occurrence by its own binding rule; a token that CANNOT be coherently re-bound (its anchor names the unperturbed version essentially) makes that perturbation inadmissible for V6 and reportable under V3c instead — silence about rebinding is not permitted. Metamorphic fixtures (paired correct-without-marker / incorrect-with-marker probes) are V6's operational estimator under clause (i). \textbf{V6 is not V2b-P:} V2b-P asks whether the configured procedure is truth-conducive over its declared reference class and criterion; V6 asks whether the episode is stable under a declared neighboring perturbation family. One test suite may supply evidence for both; the questions differ, and these definitions introduce no entailment between them.

\end{itemize}
\textbf{Result-free pathway core (Decision 0003).}


\begin{verbatim}
CorePath = { V2a, V2b-P, V2c, V3a, V3b, V3c, V3d, V3e, V4a, V5, V6 }
\end{verbatim}
Excluded by construction: \textbf{V1} (result correctness), \textbf{V2b-T} (factive token-level truth linkage — entails claim truth), and \textbf{V4b} (route near-optimality — an advisory/quality verdict recorded separately; a route can be safe and adequate without being optimal).

\textbf{Governed applicability.} Pathway applicability is DERIVED, never a discretionary post-hoc list — an executor cannot exempt a failing or untested verdict by omitting it:


\begin{verbatim}
ReqPath(e) = CorePath ∩ RequiredBy( A(e), episode-shape(e),
               risk-class(e), claims(e), governance(e) )
\end{verbatim}
Minimum derivation rules: \textbf{V3c} is required exactly when $\operatorname{MetaTok}(e) \neq \emptyset$ (the M1 zero-burden rule is preserved: $\operatorname{MetaTok}(e) = \emptyset \Rightarrow \operatorname{V3c} \notin \operatorname{ReqPath}(e)$); \textbf{V4a} is required when the episode selects or authorizes a route; \textbf{V5} is required when the episode makes a closure/completion claim or dispositions residuals; \textbf{V6} is required when $A(e)$ or the risk class predeclares a robustness obligation. \textbf{"Not tested" is not "not applicable":} every verdict excluded from $\operatorname{ReqPath}(e)$ carries a recorded applicability reason on the episode record. ($\operatorname{ReqPath}(e)$ is the sole pathway requirement function; the earlier separately recorded applicability set $\operatorname{App}(\cdot)$ is retired by Decision 0005 — nothing it expressed is lost, since exclusions now live in the per-verdict \texttt{not-applicable} statuses and reasons.)

\textbf{Machine-readable derivation and honest status (R3).} The repository ships $\operatorname{RequiredBy}$ as a machine-readable governance rule table (\texttt{docs/\allowbreak{}governance-\allowbreak{}requirements.\allowbreak{}yaml}) over a typed episode-shape input contract (\texttt{has\_meta\_tokens}, \texttt{selects\_route}, \texttt{makes\_closure\_claim}, \texttt{robustness\_obligation}, \texttt{risk\_class}), and \texttt{scripts/\allowbreak{}derive\_\allowbreak{}reqpath.\allowbreak{}py} derives $\operatorname{ReqPath}(e)$ from it with a per-verdict derivation trace (required / not-required, governing rule, rationale); fixtures in \texttt{tests/\allowbreak{}reqpath-\allowbreak{}fixtures.\allowbreak{}json} pin the derivation, including an omission-attack case in which a declared path missing a derivable requirement is detected rather than tolerated. Honest scope: the shipped table is \emph{one} complete, deterministic governance instance. $\operatorname{RequiredBy}$ in general remains a \textbf{governance-supplied parameterized interface} — the theory does not claim a closed universal calculus over every conceivable governance regime, and this part of the formalism is accordingly an open parameter, not "closed."

\textbf{Verdict status is not Boolean.}


\begin{verbatim}
Status_i(e) ∈ { pass, fail, undetermined, not-applicable }

PathwayAdequate(e)      iff every V_i ∈ ReqPath(e) has status pass
PathwayDefective(e)     iff some  V_i ∈ ReqPath(e) has status fail
PathwayUndetermined(e)  iff no required verdict fails and at least
                        one required verdict is undetermined
\end{verbatim}
A missing assessment is \texttt{undetermined} — never silently counted as pass, never silently removed as not-applicable.

\textbf{Declared implication/dependency table} (complete; nothing else is asserted, and pairwise independence is NOT claimed where unproven):


\begin{center}
\begin{tabular}{@{}p{0.470\linewidth}p{0.470\linewidth}@{}}
\toprule
\textbf{Relation} & \textbf{Status} \\
\midrule
$\operatorname{V2b-T}_q(e) \to \operatorname{V1}_q(e)$ & \textbf{Definitional entailment} (factive, claim-wise) \\
$\operatorname{V2b-P}(e) \to \operatorname{V1}(e)$ & \textbf{No} — non-factive by construction \\
$\operatorname{V3e}(e) \to \operatorname{V1}(e)$ & \textbf{No} — decision-time reasonableness is result-free \\
$\operatorname{PathwayAdequate}(e) \to \operatorname{V1}(e)$ & \textbf{No} — the core excludes V1 and V2b-T \\
$\operatorname{V1}(e) \to \operatorname{PathwayAdequate}(e)$ & \textbf{No} — stopped-clock and defective-binding cases \\
$\operatorname{V1}(e)$ at the governed level vs claim-wise $\operatorname{V1}_q(e)$ & Definitional: profile-level V1 aggregates the placed claims per the declared level \\
$\operatorname{PathwayAdequate}(e) \leftrightarrow \neg \operatorname{PathwayDefective}(e)$ & \textbf{Not asserted} — \texttt{undetermined} is a third state \\
V6 vs V2b-P & \textbf{No entailment either way} (distinct questions; shared evidence possible) \\
All other pairs & No definitional implication introduced by these definitions; empirical correlation is never entailment \\
\bottomrule
\end{tabular}
\end{center}

\subsection*{4.2 Result × pathway matrix}
The matrix applies only AFTER the pathway status is resolved as adequate or defective; $\operatorname{PathwayUndetermined}$ episodes remain outside the four resolved cells until audited. (Decision 0003: this supersedes the earlier parenthetical that committed the token-level truth-connection verdict to the conjunction — the incorrect-result/adequate-pathway cell is genuinely representable.)


% breakable-row-table: data-rows=2 total-rendered-characters=1627 max-row-rendered-characters=981
\par\medskip
\hrule height 0.8pt
\smallskip
% breakable-row: 1/2
\noindent\textbf{Column 1}: \textbf{Result correct (V1)}\par
\noindent\textbf{PathwayAdequate}: \emph{Nominal.} A current behavioral validator, correctly bound to the current artifact (V3c pass where applicable), under a sufficiently reliable configured procedure, returns the correct result. Nothing to fix.\par
\noindent\textbf{PathwayDefective}: \emph{Correct + defective (stopped-clock, compensating error, or governing-side luck).} Possible failure loci: V2b-P, V3a, V3b, V3c, V3d, or V6 — the marker validator that happens to be right today; the mis-bound metaorthemma whose wrong reference plane locally coincides with the right one. Remedy targets the procedure, rule, or binding — not the verdict.\par
\smallskip
\hrule height 0.4pt
\smallskip
% breakable-row: 2/2
\noindent\textbf{Column 1}: \textbf{Result incorrect (¬V1)}\par
\noindent\textbf{PathwayAdequate}: $\operatorname{AdequatePathError}(e) := \neg \operatorname{V1}(e) \wedge \operatorname{PathwayAdequate}(e)$ — the justified rare miss of a non-perfect but sufficiently reliable process used correctly: a genuine, representable cell, expected whenever such a process runs. Where V3e ∈ ReqPath(e), $\operatorname{JustifiedMiss}(e) := \operatorname{AdequatePathError}(e)$ — V3e's pass is already inside PathwayAdequate, so no redundant conjunct is added. Remedy, if any: a new discriminating evidence source; no rule or executor blame. \textbf{Note:} orthemic adequacy does not establish moral, legal, institutional, or theological blamelessness — "blameless" is not a formal predicate here.\par
\noindent\textbf{PathwayDefective}: \emph{Compound failure.} Wrong placement via a defective pathway — e.g., a deploy gate reads a cached test result from the previous commit and approves the current one (¬V1 with V2c and V3a failures). Remedy: both the placement and the rulebook/lineage machinery.\par
\smallskip
\hrule height 0.8pt
\par\medskip
\textbf{Worked verdict fixtures (deterministic; the four resolved cells are all satisfiable, plus the undetermined state — machine-checkable encodings in \texttt{tests/\allowbreak{}verdict-\allowbreak{}fixtures.\allowbreak{}json}, validated by \texttt{scripts/\allowbreak{}validate\_\allowbreak{}verdict\_\allowbreak{}semantics.\allowbreak{}py}):}


% breakable-row-table: data-rows=7 total-rendered-characters=2343 max-row-rendered-characters=547
\par\medskip
\hrule height 0.8pt
\smallskip
% breakable-row: 1/7
\noindent\textbf{Fixture}: \textbf{F1 — nominal}\par
\noindent\textbf{Setup}: Current behavioral validator, correctly bound (V3c pass where applicable), reliable configured procedure, correct result\par
\noindent\textbf{Expected statuses}: V1 pass; every ReqPath verdict pass ⇒ \textbf{PathwayAdequate}\par
\smallskip
\hrule height 0.4pt
\smallskip
% breakable-row: 2/7
\noindent\textbf{Fixture}: \textbf{F2 — stopped clock}\par
\noindent\textbf{Setup}: SCAN-CLEAN marker validator happens to be correct today but fails its declared reliability and perturbation criteria\par
\noindent\textbf{Expected statuses}: V1 pass; V2b-P fail; V6 fail ⇒ \textbf{PathwayDefective} (V2b-T may also fail as a result-side diagnosis: correct, but not through the truth-relevant mechanism)\par
\smallskip
\hrule height 0.4pt
\smallskip
% breakable-row: 3/7
\noindent\textbf{Fixture}: \textbf{F3 — justified rare miss}\par
\noindent\textbf{Setup}: Properly configured, faithfully executed procedure meeting its declared reference-class reliability threshold errs on this token\par
\noindent\textbf{Expected statuses}: V1 fail; V2b-P pass; V3a/V3b/V3c/V3d/V3e pass where applicable ⇒ \textbf{PathwayAdequate}; V2b-T\_q \textbf{fails} (the placed claim is false — a false claim's truth-linkage verdict fails by factivity; it is never recorded "unavailable" absent an explicit declared applicability rule). \emph{This fixture is mandatory: it proves the lower-left cell is no longer definitionally blocked.}\par
\smallskip
\hrule height 0.4pt
\smallskip
% breakable-row: 4/7
\noindent\textbf{Fixture}: \textbf{F4 — defective metaorthemma, lucky result}\par
\noindent\textbf{Setup}: Sound standard and policy, faithful execution; the concrete token binds the wrong reference plane / tolerance / fixture / success surface / calibration; the result happens to be correct\par
\noindent\textbf{Expected statuses}: V1 pass; V3a pass; V3b pass; V3d pass; \textbf{V3c fail} ⇒ \textbf{PathwayDefective}\par
\smallskip
\hrule height 0.4pt
\smallskip
% breakable-row: 5/7
\noindent\textbf{Fixture}: \textbf{F5 — unresolved audit}\par
\noindent\textbf{Setup}: The result is known, but required robustness or provenance evidence has not yet been evaluated\par
\noindent\textbf{Expected statuses}: one required verdict \texttt{undetermined} ⇒ \textbf{PathwayUndetermined}; neither PathwayAdequate nor PathwayDefective is asserted\par
\smallskip
\hrule height 0.4pt
\smallskip
% breakable-row: 6/7
\noindent\textbf{Fixture}: \textbf{F6 — stale directive} (Decision 0006; worked case in \texttt{examples/\allowbreak{}compaction-\allowbreak{}stale-\allowbreak{}steer.\allowbreak{}md})\par
\noindent\textbf{Setup}: A recovered governing directive is authentic but superseded; the executor follows it faithfully; the closure claim asserts it was in force\par
\noindent\textbf{Expected statuses}: V1 fail; V2c fail; \textbf{V3c fail} (token not Current); V3d pass; V5 fail ⇒ \textbf{PathwayDefective} — faithful execution under a defective governing token\par
\smallskip
\hrule height 0.4pt
\smallskip
% breakable-row: 7/7
\noindent\textbf{Fixture}: \textbf{F7 — safe but suboptimal route}\par
\noindent\textbf{Setup}: The chosen route is admissible but demonstrably wasteful\par
\noindent\textbf{Expected statuses}: V4a pass; \textbf{V4b fail} (advisory, outside the core) ⇒ \textbf{PathwayAdequate} — safe ≠ best, and route quality never defeats adequacy\par
\smallskip
\hrule height 0.8pt
\par\medskip

\subsection*{4.3 Epistemic constraint and nearest neighbors}
The defensible novelty claim is exactly this and no more: \textbf{the prior manuscript, and the operational discipline it describes, did not explicitly represent the concrete classification/handling episode as an auditable occurrence distinct from its result.} It is \emph{not} claimed that no established field can represent episodes. Nearest neighbors, honestly: \textbf{process reliabilism} (epistemology) names precisely the correct-result-through-unreliable- process structure and the truth-connection idea — for \emph{beliefs}; it offers no operational record schema (no typed evidence, route, residual disposition, successor occurrence, or perturbation estimator for agent/validator episodes). \textbf{Audit trails / provenance records} (W3C PROV, audit logs, ML lineage systems) are operational and token-level — who did what to which artifact when — but typically carry no candidate set, plural inferred profile, per-claim evidence status, residual-disposition ledger, or verdict layer separating result correctness from pathway adequacy and robustness. \textbf{Assurance practice} (control-effectiveness vs outcome testing) and \textbf{metamorphic testing} each contain one verdict (V3a-like, V6-like) without the joint object. What the episode signature adds is the \emph{joint} auditable object on which V1–V6 are simultaneously definable; whether that pays for itself over "audit the validator when suspicious" is an open empirical question (fixture E5), not a settled fact.


\bigskip\hrule\bigskip

\section*{5. Mechanical and distributed orthing (2E)}

\subsection*{5.1 Mechanical executors}
Nothing in the episode signature or the verdicts quantifies over awareness. Any rule-governed, evidence-updating executor — a predictive-text decoder resolving an ambiguous keypress sequence, a CI validator, a review chain — executes orthing, and its episodes bear \textbf{all applicable verdicts} ($\operatorname{ReqPath}(e)$ plus recorded not-applicable reasons, §4.1; the earlier "all six verdicts" contradicted partial applicability and is corrected). Stopped-clock analogs arise mechanically (a decoder emits the right word from a frequency prior that would misfire on the neighboring input), so V6 is meaningful without any conscious subject. \textbf{Consciousness is not required.} The actor field $\alpha$ is populated with the mechanical executor — never left empty — and executor identity is distinct from any authority index.


\subsection*{5.2 Distributed episodes: graph and composition}
Consider one case handled across sensors, validators, agents, routers, a human sign-off, and downstream consumers. Model it twice, at two levels:

\textbf{Sub-episode graph.} $\Gamma_E = (E, \rightsquigarrow)$ is a finite DAG whose nodes are episodes $e_1,\dots,e_n$ (each with the full §2 signature — own actor $\alpha_i$, own $\vec{\mu}_i$) and whose edges are typed: \textbf{handoffs} ($e_i \rightsquigarrow e_j$ iff a projection of $\operatorname{out}(e_i)$ — a partial profile, an evidence item with its scope, a route, a disposition — appears in the input or evidence $H_j$ of $e_j$; gloss: what one worker concluded is what the next starts from), \textbf{supersession} (a later episode revises an earlier one's conclusions), and \textbf{retry-of}. Handoffs are the loci of transport error: a verdict crossing an edge without its scope and version is how stale evidence propagates.

\textbf{Edge orientation, one convention.} Every edge of $\Gamma_E$ is oriented from the EARLIER node to the LATER node; edge LABELS carry the semantics. A \texttt{handoff} edge $e_i \rightsquigarrow e_j$ means $e_j$ consumes a projection of $\operatorname{out}(e_i)$; a \texttt{supersession} edge $e_i \rightsquigarrow e_j$ means the later $e_j$ revises $e_i$'s conclusions; a \texttt{retry-of} edge $e_i \rightsquigarrow e_j$ means $e_j$ re-attempts $e_i$'s task. (Prose like "a supersession edge to its predecessor" reads the label backward along the same earlier→later edge; the graph stores one orientation only.)

\textbf{Why a DAG is right despite retries and revision loops.} Nodes are \emph{dated tokens}; every typed edge points forward in time under the convention above, so the token graph is acyclic by construction. A retry or rework loop does not add a back edge — it creates a \emph{new} episode node reached by a \texttt{retry-of} or \texttt{supersession} edge from its predecessor. The genuinely cyclic structure (the same policy re-entered, review returning work to an author-role) lives at the TYPE/policy level, which the formalism represents as repeated instantiation of the same $\pi$, never as cycles among tokens.

\textbf{Composition.} The graph composes into ONE episode $e_\Gamma = \operatorname{comp}(\Gamma_E)$ — a full §2 signature at the boundary level — iff:


\begin{enumerate}
\item \textbf{One case, one analysis:} all $e_i$ address $m$ or its $\operatorname{succ}$-lineage under a single declared analysis $A$ (hence one task $T = \operatorname{task}(A)$; two sub-episodes sharing a task but differing in tolerance, representation, or boundary do NOT compose without an explicitly declared fusion analysis);

\item \textbf{A declared boundary:} governance names the composite actor $\alpha_{e_\Gamma}$ (pipeline, team, institution) and configuration $\vec{\mu}_{e_\Gamma}$, including precedence across sub-episode rules;

\item \textbf{A declared fusion rule:} $\hat{p}_{e_\Gamma}$, $\operatorname{estatus}_{e_\Gamma}$, $\delta_{e_\Gamma}$ are a stated aggregation of sub-outputs (e.g., profile union with per-claim status = weakest status along its supporting path; residual ledger = merged ledgers with surviving ownership) — well-defined, not read off the last node.

\end{enumerate}
Then $e_\Gamma$ is a token in its own right, and each $e_i$ becomes an orthemma of the composite-level audit via the reification embedding ($\iota_n(e_i) \in M^{(n+1)}$ per §1 — never literal membership). \textbf{Answer to one-vs-graph: both, at different levels.} The graph is the true description at the component level; the composed episode at the boundary level; $\operatorname{comp}$ and its three conditions say when the second description exists. When condition 2 or 3 fails — no declared boundary, no fusion rule — there is only the graph, and talk of "the system decided" is unwarranted composite attribution.

Composite verdicts do not reduce to conjunctions: $\operatorname{correct}(e_\Gamma)$ can hold while some $\operatorname{correct}(e_i)$ fails (a downstream validator caught it), and every $\operatorname{correct}(e_i)$ can hold while $\operatorname{correct}(e_\Gamma)$ fails (each verdict true in scope, the composition transporting one out of scope). Hence the composite level must be audited separately.


\subsection*{5.3 Multi-actor episodes over one occurrence lineage (derived)}
Several actors' episodes can address ONE shared occurrence lineage with analysis-indexed, possibly opposed evaluations (competitive games are the clean case; the full stress test and category corrections live in \texttt{orthemic-\allowbreak{}multi-\allowbreak{}actor-\allowbreak{}conflict-\allowbreak{}note.\allowbreak{}md}). This is a \textbf{multi-analysis context}: each actor $\alpha$ evaluates under its own declared analysis $A_\alpha$ (shared frame components, diverging at least in task $T_\alpha = \operatorname{task}(A_\alpha)$), the task-relative abbreviation is forbidden, and ground truth is written in full, $O^*(m; A_\alpha)$. Actor indices are written $\alpha, \beta$ here — the letter $A$ is reserved for the analysis. Derived definitions — clarifying extensions of the existing actor/analysis indices, \textbf{not a new formal addition}:


\begin{itemize}
\item $x_\alpha = \Omega_\alpha(m)$ — actor-α's OBSERVATION of the shared occurrence (imperfect information is divergent observation over ONE occurrence: $\Omega_\alpha(m) \neq \Omega_\beta(m)$; perfect information is the special case of full, equal observation) — the observation/occurrence distinction applied per actor;

\item $\hat{p}(e, \alpha, T_\alpha)$ — actor-α's current inferred profile in episode $e$ under its analysis $A_\alpha$ (task $T_\alpha$), inferred FROM $x_\alpha$, never from $m$ directly;

\item $\operatorname{GoalSchema}(\cdot)$ — a PARAMETRIC ortheme schema (the shared goal-form under role substitution); grounded instantiations $\operatorname{GoalSchema}(\alpha)$, $\operatorname{GoalSchema}(\beta)$ are DISTINCT targets even when the schema is one (schema identity ≠ grounded-ortheme identity ≠ target-profile overlap);

\item $\mathcal{G}_{\alpha,A_\alpha} \subseteq \Pi_{A_\alpha}$ — actor-α's grounded target profile set — NORMATIVE TYPING: a SET of profiles, each member one complete profile (the task's appropriateness rendered as a profile set; targets are selected BY the task, which remains an external parameter — a target profile is not an objective and an objective is not an ortheme);

\item $\phi(\alpha \to \beta)$ — a role/perspective isomorphism between $\Pi_{A_\alpha}$ and $\Pi_{A_\beta}$ where one exists (goal-schema invariance with disjoint extensions: $\mathcal{G}_{\beta,A_\beta} = \phi(\mathcal{G}_{\alpha,A_\alpha})$, possibly $\mathcal{G}_\alpha \cap \mathcal{G}_\beta = \emptyset$);

\item $\operatorname{Conflict}_m(\mathcal{G}_\alpha, \mathcal{G}_\beta)$ / $\operatorname{Compat}_m(\mathcal{G}_\alpha, \mathcal{G}_\beta)$ — WELL-TYPED across profile spaces: the two target sets live in different spaces ($\mathcal{G}_\alpha \subseteq \Pi_{A_\alpha}$, $\mathcal{G}_\beta \subseteq \Pi_{A_\beta}$), so they are never compared by bare set intersection. With $\operatorname{Reach}(m)$ the occurrences reachable from $m$ in the successor structure,

\begin{verbatim}
Compat_m(𝒢_α, 𝒢_β)   iff  ∃ m′ ∈ Reach(m):
                            O*(m′; A_α) ∈ 𝒢_α  ∧  O*(m′; A_β) ∈ 𝒢_β
Conflict_m(𝒢_α, 𝒢_β) iff  no such reachable m′ exists
\end{verbatim}
— one shared occurrence, evaluated under each analysis separately. Cooperation is joint realizability with a shared analysis and shared target set ($A_\alpha = A_\beta$, $\mathcal{G}_\alpha = \mathcal{G}_\beta$), or an explicit alignment map $\phi(\alpha \to \beta)$ identifying targets across spaces — never bare equality of sets inhabiting different profile spaces. Zero-sum games also admit draw successors realizing neither target set.

\end{itemize}
Coupled episodes are the §5.2 graph over the shared lineage: each actor's action creates the successor occurrence the other's episode consumes. Game transition rules are DOMAIN structure, utility is the TASK, minimax is a POLICY — none is a metaortheme. The LIMITED sense in which competitors share a metaorthemic configuration: both may consult the same governing distinctions (terminality, evidence-validity about the opponent's information state, analysis-sufficiency) while holding opposed targets — configurations govern HOW placements are made, never WHOSE target wins. Occurrence identity in games includes position, history, and player-to-move — worldly facts (\#4 applied); a player's INFORMATION SET is NOT identity but that player's OBSERVATION $x_\alpha = \Omega_\alpha(m)$ (a correction adopted after external review — filing it into $\kappa$ would invert the occurrence/observation line). Not all orthemes are actor-relative — actor-relativity enters exactly where the task does.


\subsection*{5.4 Semantic depth as an orthogonal axis}
Define, informally, $\operatorname{depth}(e)$ as the executor's grasp of the governing distinctions: hard-coded check < model-based inference over declared alternatives < reflective capacity to propose revisions $\rho(\mu, \mathrm{lesson})$ itself. Depth is orthogonal to every §4 verdict — shallow executors can run adequate pathways and deep ones defective pathways. Depth matters operationally in one place only: which revisions the executor may perform locally versus must escalate across the governance boundary. It is an axis of the executor, not of the episode's correctness or adequacy.


\bigskip\hrule\bigskip
\textbf{Optional latent-variable extension (Decision 0015).} A declared analysis $A$ may optionally carry sequential latent-variable apparatus ($Z_A$, transitions $P_A$, latent candidate sets/posteriors, and concrete internal representations $y$). That apparatus is \textbf{non-primitive and optional}: it adds nothing to the settled ontology of this formalization, a latent state is not an ortheme, a posterior is not $O^*(m; A)$, and representational geometry is neither necessary nor sufficient for orthemic distinctness or for pathway adequacy. The bridge $\operatorname{ProfileOf}_A \subseteq Z_A \times \Pi_A$ is partial and must be exhibited, never assumed; latent labels are non-identifiable up to relabelling absent declared anchoring or a validated alignment map. \textbf{Every core claim of this formalization remains statable without the latent layer.} Declaring it is a version event on $A$. See \texttt{docs/\allowbreak{}decisions/\allowbreak{}0015-\allowbreak{}latent-\allowbreak{}state-\allowbreak{}observation-\allowbreak{}and-\allowbreak{}representation-\allowbreak{}boundary.\allowbreak{}md} and \texttt{docs/\allowbreak{}related-\allowbreak{}work/\allowbreak{}LATENT-\allowbreak{}STATE-\allowbreak{}INFERENCE-\allowbreak{}AND-\allowbreak{}ORTHEMOLOGY.\allowbreak{}md}.


\section*{Status ledger}
\textbf{Accepted (analytic/definitional):} level-indexing via $\iota_n$ with no episode-referent "meta-orthemma" noun (§1, counterargument recorded; per the §1 qualification, the Decision-0002 configuration token $\bar{\mu}$ is a different object, carried in §2.2/§3/§4.1 V3c); the record-style episode signature with typed candidate families, claim ledger, handoffs, warrant, and labeled successors (§2.2); endpoint-plus-bounded-trace with governance-declared granularity (§2.2a); split model (C) with the $\{\mathrm{O}, \mathrm{I}, \mathrm{E}, \mathrm{D}, \mathrm{R}, \mathrm{V}, \mathrm{W}\}$ boundary (§3, downsides preserved); verdict separability with governance-derived $\operatorname{ReqPath}(e)$ (§4.1); token-level DAG with typed edges (§5.2). \textbf{Rejected:} literal set inclusion for episodes; forced single successor; blanket edition-wide revision reopening; safeguard (v) word-utility test inside concept admission. \textbf{Unresolved alternatives:} safety/reliabilist/calibration variants of V2b-P remain admissible declared criteria; single-object metaortheme presentation (C1) remains permitted as user-facing shorthand. \textbf{Experimental gates:} episode reification's practical delta over "audit the validator when suspicious" — fixture E5; vocabulary utility of every coined term — the Arm A/B/C benchmark (designed, NOT run). \textbf{Evidence tier:} machine-checked internal agreement over the declared definitions, schemas, examples, and adversarial fixtures — not a proof of mathematical consistency or completeness; no public observational dataset exists; operational usefulness is untested; nothing here is experimentally validated. \textbf{Provenance:} \texttt{docs/\allowbreak{}provenance/\allowbreak{}document-\allowbreak{}history.\allowbreak{}md}; this file is the canonical integrated body.


\part*{Multi-Actor / Zero-Sum Orthemic Note — stress test, corrections, and formal residue}
\textbf{Status (revision R6, 2026-07-20; fresh-session repository review completed; not external human peer review; not empirically validated):} derived multi-actor extension of the formal core; typing follows decisions 0009–0015; nothing here is experimentally validated.


\begin{quote}
\textbf{Provenance.} Written from three informal zero-sum discussion notes supplied to the project (private; not published), treated as a THEORY STRESS TEST, not accepted doctrine — original header preserved in \href{\detokenize{../docs/provenance/document-history.md}}{\texttt{docs/\allowbreak{}provenance/\allowbreak{}document-\allowbreak{}history.\allowbreak{}md}}. Verdict up front: the source notes surface a real and useful multi-actor layer, AND they conflate five distinct categories; the repair below preserves the insight without the conflations. The formal residue is a \textbf{clarification plus derived extensions of the existing actor/analysis indices — not a seventh formal addition.}

\end{quote}

\section*{1. What the source notes get right (preserved)}

\begin{enumerate}
\item \textbf{One shared concrete occurrence.} Both players face the SAME orthemma — the game state. Multi-actor cases do not multiply occurrences; they multiply perspectives on one occurrence lineage.

\item \textbf{Actor/analysis-relative evaluation can oppose.} "Correct for A" and "correct for B" genuinely diverge under a competitive task; the theory's per-actor analysis indexing ($A_\alpha$ with $T_\alpha = \operatorname{task}(A_\alpha)$, $\Pi_{A_\alpha}$) was built for exactly this and handles it. (Notation update, D1 reconciliation: $A$ is reserved for the declared analysis; the formal actor indices are $\alpha, \beta$, with the prose labels "Player A / Player B" naming the actors $\alpha, \beta$ respectively. Multi-actor evaluation is a multi-analysis context, so the task-relative shorthand of manuscript Definition 3 is forbidden here and ground truth is written in full, $O^*(m; A_\alpha)$.)

\item \textbf{Structural isomorphism across actors.} A's and B's reasoning can be the same schema under role substitution while their targets are materially incompatible — a real and clarifying observation.

\item \textbf{Coupled episodes through a shared successor.} A's move-episode produces the successor occurrence B's episode consumes: the typed-edge episode graph (core-formalization §5.2) over one occurrence lineage.

\item \textbf{Competitive vs cooperative as target-set relations.} Pair programming = shared target profile; chess = disjoint target profiles. Correct.

\end{enumerate}

\section*{2. Category conflations corrected}
\textbf{C1 — "Player A's ortheme: the state where my utility is maximized."} That is not an ortheme; it is an OBJECTIVE rendered as a state description. An ortheme is a repeatable state-type of the occurrence ("White checkmates Black" IS one). The repair: player A's \emph{target profile set} $\mathcal{G}_{\alpha,A_\alpha} \subseteq \Pi_{A_\alpha}$ is the set of terminal profiles that satisfy that player's task — the states are orthemes; "maximize my utility" is the task that SELECTS them. Utility/objective/loss are external parameters (core-formalization §3), never orthemes and never metaorthemes.

\textbf{C2 — "their orthemes are completely identical in structure."} Isomorphism is not identity. The goal SCHEMA "win-for-me" is invariant under the role substitution φ(α→β); the EXTENSIONS differ: φ maps "White checkmates Black" to "Black checkmates White." Two consequences the source notes blur: (i) nothing is "the same ortheme" here — the schemas are isomorphic, the target profiles disjoint; (ii) the interesting fact is precisely that φ exists and 𝒢\_\{β,A\_β\} = φ(𝒢\_\{α,A\_α\}) while 𝒢\_\{α,A\_α\} ∩ 𝒢\_\{β,A\_β\} = ∅.

\textbf{C3 — "the metaortheme = the rules of the game + the logic of minimax."} Three different things, none of them a metaortheme:


\begin{itemize}
\item \textbf{game transition rules} — the DOMAIN's lawful-transition structure (which successor occurrences are reachable at all);

\item \textbf{utility / rational maximization} — the TASK/objective, an external parameter excluded from governed components;

\item \textbf{minimax} — a POLICY/strategy (a way of choosing actions under the task). A metaorthemic configuration in a game would govern something like: what counts as EVIDENCE about the opponent's information state; which analysis depth suffices before a move may be made (an evidence-sufficiency distinction); whether an engine evaluation is validated or provisional; when a position assessment is stale (position edited between analyses). The source notes' "neutral governing logic" intuition is real, but the neutral shared things are the domain rules and the task-form — the metaortheme slot was mislabeled.

\end{itemize}
\textbf{C4 — occurrence identity is thinner than it should be, and information sets are OBSERVATIONS, not identity (corrected after external review).} "The board" is not enough: chess needs position + move history (castling/repetition rights) + player-to-move — all WORLDLY facts of the one occurrence, so they belong in the identity key κ (\#4 applied). But a player's INFORMATION SET is not part of κ: in poker the one concrete deal (all hole cards) IS the occurrence m; what player α knows of it is α's OBSERVATION, x\_α = Ω\_α(m) — the framework's own observation-vs-occurrence distinction, applied per actor. Imperfect information = one shared m with divergent Ω\_α(m) ≠ Ω\_β(m); perfect-information games are the special case Ω\_α = Ω\_β = full state. Filing the information set into κ would invert the occurrence/observation line the theory is built on.

\textbf{C5 — strict zero-sum does not mean every state realizes ±X.} U\_α + U\_β = 0 admits draws (0,0). Terminal profiles partition into 𝒢\_α, 𝒢\_β, and a nonempty draw set in general; "the objective manifestation of A's ortheme mathematically prevents B's" holds only between the two win-sets, not over the whole terminal space.

\textbf{C6 — "reflexively selfish" is task language, not psychology or ethics.} Under T\_α = "win for the player labeled A," correctness is α-indexed by construction; that is a property of the DECLARED task, not a discovery about agents. And not all orthemes are actor-relative: "the file hash matches," "the parse is grammatical" are actor-independent placements; actor-relativity enters exactly where the task does.


\section*{3. Formal residue (assessed against the existing core)}
Already expressible with existing machinery: episodes carry actor α and declared analysis A\_α (task T\_α = task(A\_α)); placements are actor/analysis-indexed; the episode graph couples episodes over one lineage. \textbf{New derived definitions worth adding (clarification, NOT a new formal addition — the six-addition count is unchanged):}


\begin{itemize}
\item $x_\alpha = \Omega_\alpha(m)$ — actor-$\alpha$'s \textbf{observation} of the shared occurrence $m$ (imperfect information: $\Omega_\alpha(m) \neq \Omega_\beta(m)$ over one $m$);

\item $\hat{p}(e, \alpha, T_\alpha)$ — the current inferred profile of actor $\alpha$'s episode $e$ under its analysis $A_\alpha$, task $T_\alpha = \operatorname{task}(A_\alpha)$ (inferred from $x_\alpha$, never from $m$ directly);

\item $\operatorname{GoalSchema}(\cdot)$ — a \textbf{parametric} ortheme schema, the shared goal-form under role substitution (win-for-$\alpha$);

\item $\mathcal{G}_{\alpha,A_\alpha} \subseteq \Pi_{A_\alpha}$ — actor-$\alpha$'s grounded \textbf{target profile set}: the profiles $T_\alpha$ declares appropriate as outcomes — the schema's instantiation at $\alpha$. Normative typing: a \emph{set} of profiles, each member one complete profile;

\item $\phi(\alpha \to \beta)$ — a role/perspective isomorphism between $\Pi_{A_\alpha}$ and $\Pi_{A_\beta}$, when one exists (schema invariance: $\mathcal{G}_{\beta,A_\beta} = \phi(\mathcal{G}_{\alpha,A_\alpha})$);

\item $\operatorname{Compat}_m(\mathcal{G}_\alpha, \mathcal{G}_\beta)$ — $\exists\, m' \in\allowbreak{} \operatorname{Reach}(m):\ O^*(m'; A_\alpha) \in\allowbreak{} \mathcal{G}_\alpha \wedge\allowbreak{} O^*(m'; A_\beta) \in\allowbreak{} \mathcal{G}_\beta$ (well-typed: one shared occurrence, each analysis evaluating in its own profile space — never bare intersection of sets from different spaces);

\item $\operatorname{Conflict}_m(\mathcal{G}_\alpha, \mathcal{G}_\beta)$ — no such reachable $m'$ exists (cooperation: shared analysis $+$ shared target set, or an explicit alignment map $\phi(\alpha \to \beta)$).

\end{itemize}
\textbf{Three distinct relations the parametric schema untangles:} (i) SCHEMA identity — players A and B run the same parametric form $\operatorname{GoalSchema}(\cdot)$; this is the true content of the source notes' "locally identical" intuition; (ii) GROUNDED-ORTHEME identity — $\operatorname{GoalSchema}(\alpha)$ and $\operatorname{GoalSchema}(\beta)$ are DISTINCT grounded targets, never identical in a strictly competitive task; (iii) TARGET-PROFILE overlap — $\mathcal{G}_\alpha \cap \mathcal{G}_\beta$ may be empty (zero-sum win-sets), partial, or total (cooperation). Local similarity is (i); material exclusivity is emptiness at (iii); no level asserts identity of the players' orthemes.

\textbf{The limited sense in which competitors DO share a metaorthemic configuration:} both players may consult the SAME governing distinctions — terminality (decided vs live position), evidence-validity for reading the opponent's information state, analysis-sufficiency before moving. Sharing those configurations while holding opposed targets is exactly the configuration-vs-task separation: metaorthemes govern HOW placements are made and validated, never WHOSE target wins.

Joint-episode picture: A's and B's orthing episodes alternate on one occurrence lineage; each actor's action creates the successor the other's episode consumes (typed handoff edges); under Conflict, each episode's route choice is also an attempt to steer the lineage away from the other's target set. Draw states are successors in neither target set.


\section*{4. Destinations}

\begin{itemize}
\item \textbf{Core formalization:} §5.3 there (multi-actor episodes over one lineage; the definitions above; explicitly derived, not a new addition). Integrated.

\item \textbf{Terminology benchmark:} a multi-actor fixture — a scenario where the arm must keep apart: shared occurrence / actor-relative target / current profile / objective / policy / metaorthemic distinction. The source notes become the NEGATIVE calibration example for objective-vs-metaortheme-vs-policy conflation (joining the standing objective-mislabeled-as-metaortheme example). Integrated into the designed (unrun) protocol.

\item \textbf{Revised manuscript:} §10 worked example (chess) with C1–C6 corrections — present in \texttt{orthemma-\allowbreak{}ortheme-\allowbreak{}systems-\allowbreak{}revised-\allowbreak{}draft.\allowbreak{}md}.

\item \textbf{Operational product lane:} \textbf{no new issue.} After ordinary-language translation ("two agents with opposed success criteria acting on one artifact") no product invariant emerges that identity, success-surface, and handoff governance do not already cover; competitive multi-agent execution is out of product scope. Default expectation confirmed: theory + benchmark only.

\end{itemize}

\section*{5. Status}
Analytic note; no empirical claim. The multi-actor fixture is part of the UNRUN terminology benchmark. Nothing here alters the six formal additions, the verdict vector, or any product issue.


\twocolumn
\bibliographystyle{plainnat}
\bibliography{../../../references/orthemology}
\end{document}
